\begin{proposition}
\label{proposition-compute-bivariant}
\begin{reference}
\cite[Proposition 17.4.2]{F}
\end{reference}
Let $(S, \delta)$ be as in Situation \ref{situation-setup}.
Let $f : X \to Y$ and $g : Y \to Z$ be morphisms of schemes locally
of finite type over $S$. If $g$ is smooth of relative dimension $d$, then
$A^p(X \to Y) = A^{p - d}(X \to Z)$.
\end{proposition}

\begin{proof}
We will use that smooth morphisms are local complete intersection
morphisms whose Gysin maps exist (see Section \ref{section-koszul}).
In particular we have $g^! \in A^{-d}(Y \to Z)$. Then we can send
$c \in A^p(X \to Y)$ to $c \circ g^! \in A^{p - d}(X \to Z)$.

\medskip\noindent
Conversely, let $c' \in A^{p - d}(X \to Z)$. Denote $res(c')$ the restriction
(Remark \ref{remark-restriction-bivariant}) of $c'$ by the morphism $Y \to Z$.
Since the diagram
$$
\xymatrix{
X \times_Z Y \ar[r]_{\text{pr}_2} \ar[d]_{\text{pr}_1} & Y \ar[d]^g \\
X \ar[r]^f & Z
}
$$
is cartesian we find $res(c') \in A^{p - d}(X \times_Z Y \to Y)$.
Let $\Delta : Y \to Y \times_Z Y$ be the diagonal and denote
$res(\Delta^!)$ the restriction of $\Delta^!$
to $X \times_Z Y$ by the morphism $X \times_Z Y \to Y \times_Z Y$.
Since the diagram
$$
\xymatrix{
X \ar[r] \ar[d] & X \times_Z Y \ar[d] \\
Y \ar[r]^-\Delta & Y \times_Z Y
}
$$
is cartesian we see that $res(\Delta^!) \in A^d(X \to X \times_Z Y)$.
Combining these two restrictions we obtain
$$
res(\Delta^!) \circ res(c') \in A^p(X \to Y)
$$
Thus we have produced maps $A^p(X \to Y) \to A^{p - d}(X \to Z)$
and $A^{p - d}(X \to Z) \to A^p(X \to Y)$. To finish the proof we
will show these maps are mutually inverse.

\medskip\noindent
Let us start with $c \in A^p(X \to Y)$. Consider the diagram
$$
\xymatrix{
X \ar[d] \ar[r] & Y \ar[d] \\
X \times_Z Y \ar[r] \ar[d]^{\text{pr}_1} &
Y \times_Z Y \ar[r]_{p_2} \ar[d]^{p_1} &
Y \ar[d]^g \\
X \ar[r]^f &
Y \ar[r]^g &
Z
}
$$
whose squares are carteisan. The lower two square of this diagram
show that $res(c \circ g^!) = res(c) \cap p_2^!$ where in this formula
$res(c)$ means the restriction of $c$ via $p_1$. Looking at the upper
square of the diagram and using Lemma \ref{lemma-lci-gysin-commutes}
we get $c \circ \Delta^! = res(\Delta^!) \circ res(c)$.
We compute
\begin{align*}
res(\Delta^!) \circ res(c \circ g^!)
& =
res(\Delta^!) \circ res(c) \circ p_2^! \\
& =
c \circ \Delta^! \circ p_2^! \\
& =
c
\end{align*}
The final equality by Lemma \ref{lemma-diagonal-identity}.

\medskip\noindent
Conversely, let us start with $c' \in A^{p - d}(X \to Z)$. Looking
at the lower rectangle of the diagram above we find
$res(c') \circ g^! = \text{pr}_1^! \circ c'$.
We compute
\begin{align*}
res(\Delta^!) \circ res(c') \circ g^!
& =
res(\Delta^!) \circ \text{pr}_1^! \circ c' \\
& =
c'
\end{align*}
The final equality holds because the left two squares of
the diagram show that
$\text{id} = res(\Delta^! \circ p_1^!) = res(\Delta^!) \circ \text{pr}_1^!$.
This finishes the proof.
\end{proof}